\documentclass[11pt]{article}

\usepackage[margin=1in]{geometry}
\usepackage{graphicx}
\usepackage[colorlinks=true,allcolors=blue]{hyperref}
\usepackage{authblk}
\usepackage[style=numeric, sorting=none, backend=biber]{biblatex}
\addbibresource{references.bib}



% determines the spacing and general formatting of each section in the document
\makeatletter
\renewcommand{\@maketitle}{%
    \newpage
    \null
    \vskip 2em%
    \let\footnote\thanks
    \begin{center}%
        {\LARGE \@title \par}%
    \end{center}%
        \vskip 1em%
        {\large \noindent\begin{minipage}{\textwidth}\raggedright \@author\end{minipage}\par}%
        \vskip 1em%
    \begin{center}%
        {\@date}%
    \end{center}%
    \par
    \vskip 1.5em%
}

% Compact bibliography format: [N] Author et al., Year, ShortJournal.
\DeclareNameFormat{familyonly}{%
  \ifnumequal{\value{listcount}}{1}{%
    \namepartfamily%
    \ifnumgreater{\value{listtotal}}{1}{ et al.}{}%
  }{}%
}
\DeclareBibliographyDriver{article}{%
  \printnames[familyonly]{author}%
  , \printfield{year}%
  , \printfield{shortjournal}%
  }
\DeclareBibliographyDriver{inproceedings}{%
  \printnames[familyonly]{author}%
  , \printfield{year}%
  , \printfield{shortjournal}%
}
\AtBeginBibliography{\setlength{\parskip}{0pt}}
\setlength{\bibitemsep}{0pt}
\setlength{\bibparsep}{0pt}

\makeatother

\title{The Effects of an Icy Porous Layer on the Two-Way Radar Attenuation on Enceladus}
\renewcommand\Authands{, }
\author[1,2]{William Byrne}
\author[1]{Ana-Catalina Plesa}
\author[1]{Tina Ruckriemen-Bez}
\author[3]{Andreas Benedikter}
\author[1]{Hauke Hussmann}


\affil[1]{Institute of Space Research, German Aerospace Center (DLR), Berlin, Germany\\
\href{mailto:William.Byrne@dlr.de}{William.Byrne@dlr.de}; \href{mailto:Ana.Plesa@dlr.de}{Ana.Plesa@dlr.de}; \href{mailto:tina.rueckriemen-bez@dlr.de}{tina.rueckriemen-bez@dlr.de}; \href{mailto:Hauke.Hussmann@dlr.de}{Hauke.Hussmann@dlr.de};}
\affil[2]{University of Maryland, College Park, MD, USA\\
\href{mailto:wbyrne@umd.edu}{wbyrne@umd.edu}}
\affil[3]{Microwaves and Radar Institute, German Aerospace Center (DLR), Weßling, Germany\\
\href{mailto:Andreas.Benedikter@dlr.de}{Andreas.Benedikter@dlr.de}}
\date{01.12.2024:
\href{https://meetingorganizer.copernicus.org/EGU24/EGU24-6159.html}{View abstract on the EGU2024 website}}

\begin{document}

\maketitle

\raggedright
\setlength{\parindent}{0pt}
\setlength{\parskip}{\baselineskip}

\section*{Abstract}

Saturn's moon, Enceladus is considered a priority target for future planetary missions due to its high astrobiological potential \parencite{choblet2021enceladus}. Water jets presumably originating from a subsurface ocean have been observed at the south pole of Enceladus by NASA’s Cassini mission \parencite{hansen2006enceladus}, and their analysis provides a direct window into the ocean composition \parencite{postberg2009sodium} that, in turn, can help to understand the nature and amount of impurities that may exist within the ice shell.

Enceladus’ jet activity generates a highly porous material that affects the thermal state of the ice shell. The thickness of that layer and its distribution are poorly constrained, but local thicknesses of up to 700m have been reported from the analysis of pit chains on the surface of Enceladus \parencite{martin2023}. Such a thick porous layer can strongly attenuate the signal of radar sounders that have been proposed to investigate the Enceladus’ subsurface \parencite{soucek2023}.

Here, we use numerical simulations to determine the effects of a porous layer on the two-way radar attenuation. We generate a variety of steady-state one-dimensional thermal models based on proposed parameters for Enceladus’ ice shell thickness (5 - 35 km, \parencite{hemingway2019enceladus}), porous layer thickness (0 - 700 m \parencite{martin2023}) and its thermal conductivity (0.1 - 0.001 W/mK \parencite{seiferlin1996,ferrari2021}). In addition to systematically testing parameter combinations, we use two ice shell thickness maps \parencite{hemingway2019enceladus} together with local thermal profiles to provide a global spatial distribution of potential penetration depths that could be achieved by radar measurements. We use two material models ("high" and "low" loss) to identify the impact of chemical impurities on attenuation \parencite{kalousova17}. While the “low” loss scenario considers an ice shell composed of pure water ice, the “high” loss case is characterized by a homogeneous mixture of water ice and chlorides in concentrations extrapolated from the particle composition of Enceladus’ plume \parencite{soucek2023}.

Our results show that the presence of a porous layer has a first-order effect on the two-way radar attenuation. For regions covered by porous layers with thicknesses larger than 250 m and a thermal conductivity lower than 0.025 W/(mK) the two-way radar attenuation reaches a threshold value of 100 dB before reaching the ice-ocean interface in the low loss scenarios. In the high loss cases, for similar porous layer thicknesses and thermal conductivity, the two-way attenuation remains below 100 dB for at most 48\% of the ice shell. Depending on the local ice shell thickness and properties of the snow deposits, as little as a few percent of the ice shell can be penetrated before the 100 dB limit is reached. We note, however, that the presence of a porous layer leads to high subsurface temperatures and promotes the formation of brines at shallow depth that can be detected by future radar measurements.

\printbibliography[title={References}]

\end{document}
